\ifx \aspectratio \undefined
	\def \aspectratio{43}	% also: 169
\fi
\documentclass[xetex,mathserif,serif,aspectratio=\aspectratio]{beamer}

\usepackage{listings}
\usepackage{tikz}
\usepackage{hyperref}
\usepackage{xcolor}
\usepackage{verbatim}
\usepackage{eso-pic}
\usepackage{courier}
\usepackage{textcomp}
\usepackage[inkscapepath=/tmp/TalkSvg/]{svg}
\usepackage[cm-default]{fontspec}
\usepackage{xeCJK}

\setCJKmainfont{Source Han Sans CN}
\setCJKmonofont{Source Han Sans CN}
\setmainfont{DejaVu Sans}
\setmonofont{Courier 10 Pitch}

\hypersetup{colorlinks,linkcolor=,urlcolor=blue}

\newcommand{\link}[1]{\href{#1}{#1}}

\newcommand{\insertGraph}[3]{

	\centerline{\includegraphics[scale=#1]{#2}} % 0/0

	\centerline{#3}

}

\newcommand{\HCCLogoSimp}{
	\begin{tikzpicture}[scale=0.1,
		rev/.style={insert path={circle[radius=9]}},eo/.code={\pgfseteorule}]
		\fill [color={rgb,255:red,0; green,204; blue,255}]
	(55.371,-11.075) --
	(57.286,-11.417) --
	(55.693,-20.323) --
	(53.803,-19.976);

\begin{scope}
% \fill [color={rgb,255:red,255; green,0; blue,0}, fill opacity=0.3] (47.042,-11.705) rectangle (62.889,-27.553); 
% \fill [color={rgb,255:red,255; green,0; blue,0}, fill opacity=0.3] (54.9655,-19.629) circle [radius=7.924]; 
% \fill [color={rgb,255:red,255; green,0; blue,0}, fill opacity=0.3] (49.018,-13.682) rectangle (60.913,-25.576); 
% \fill [color={rgb,255:red,255; green,0; blue,0}, fill opacity=0.3] (54.9655,-19.629) circle [radius=5.9475]; 
% \fill [color={rgb,255:red,255; green,255; blue,0}, fill opacity=0.3] (59.4605,-11.563) -- (53.3695,-11.511) -- (54.9665,-19.629);
\clip
	(59.4605,-11.563) -- 
	(54.9665,-19.629) -- 
	(53.3695,-11.511) -- 
	(47.041,-11.511) -- 
	(47.041,-27.553) -- 
	(62.890,-27.553) -- 
	(62.890,-11.511);
\fill [color={rgb,255:red,0; green,204; blue,255}, even odd rule]
      (54.9655,-19.629) circle [radius=5.9475]
      (59.4605,-11.563) -- (53.3695,-11.511) -- (54.9665,-19.629)
      (54.9655,-19.629) circle [radius=7.924]; 
\end{scope}


	\end{tikzpicture}
}

\newcommand{\HCCLogoFull}{
	\begin{tikzpicture}[scale=0.05,
		rev/.style={insert path={circle[radius=9]}},eo/.code={\pgfseteorule}]
		\fill [color={rgb,255:red,0; green,204; blue,255}]
	(55.371,-11.075) --
	(57.286,-11.417) --
	(55.693,-20.323) --
	(53.803,-19.976);

\begin{scope}
% \fill [color={rgb,255:red,255; green,0; blue,0}, fill opacity=0.3] (47.042,-11.705) rectangle (62.889,-27.553); 
% \fill [color={rgb,255:red,255; green,0; blue,0}, fill opacity=0.3] (54.9655,-19.629) circle [radius=7.924]; 
% \fill [color={rgb,255:red,255; green,0; blue,0}, fill opacity=0.3] (49.018,-13.682) rectangle (60.913,-25.576); 
% \fill [color={rgb,255:red,255; green,0; blue,0}, fill opacity=0.3] (54.9655,-19.629) circle [radius=5.9475]; 
% \fill [color={rgb,255:red,255; green,255; blue,0}, fill opacity=0.3] (59.4605,-11.563) -- (53.3695,-11.511) -- (54.9665,-19.629);
\clip
	(59.4605,-11.563) -- 
	(54.9665,-19.629) -- 
	(53.3695,-11.511) -- 
	(47.041,-11.511) -- 
	(47.041,-27.553) -- 
	(62.890,-27.553) -- 
	(62.890,-11.511);
\fill [color={rgb,255:red,0; green,204; blue,255}, even odd rule]
      (54.9655,-19.629) circle [radius=5.9475]
      (59.4605,-11.563) -- (53.3695,-11.511) -- (54.9665,-19.629)
      (54.9655,-19.629) circle [radius=7.924]; 
\end{scope}


		\filldraw [color={rgb,255:red,0; green,170; blue,212}]
	(49.000,-46.000) rectangle (64.000,-50.500);
\filldraw [color={rgb,255:red,0; green,136; blue,170}]
	(44.500,-35.500) rectangle (49.000,-50.500);
\filldraw [color={rgb,255:red,0; green,102; blue,128}]
	(44.500,-31.000) rectangle (64.000,-35.500);
\filldraw [color={rgb,255:red,0; green,170; blue,212}]
	(27.500,-46.000) rectangle (42.500,-50.500);
\filldraw [color={rgb,255:red,0; green,136; blue,170}]
	(22.500,-35.500) rectangle (27.500,-50.500);
\filldraw [color={rgb,255:red,0; green,102; blue,128}]
	(22.500,-31.000) rectangle (42.000,-35.500);
\filldraw [color={rgb,255:red,0; green,136; blue,170}]
	(16.000,-35.500) rectangle (20.500,-51.000);
\filldraw [color={rgb,255:red,0; green,102; blue,128}]
	(4.500,-31.000) rectangle (20.500,-35.500);
\filldraw [color={rgb,255:red,0; green,68; blue,85}]
	(0.000,-24.000) rectangle (4.500,-53.500);

\begin{comment}
import re
import numpy as np
from bs4 import BeautifulSoup

def p(x) :
	return '(%.3f,%.3f)' % (x[0], -x[1])

def f(x) :
	X = x.split()
	prev_cmd = None
	i = 0
	cur = 0, 0
	def read() :
		nonlocal i, cur
		dat = X[i]; i += 1
		d1, d2 = dat.split(',')
		return cur[0] + float(d1), cur[1] + float(d2)
	def read_cur() :
		nonlocal i, cur
		dat = X[i]; i += 1
		d1, d2 = dat.split(',')
		cur = cur[0] + float(d1), cur[1] + float(d2)
		return cur
	def add(x, y) :
		return x[0] + y[0], x[1] + y[1]
	while i < len(X) :
		if len(X[i]) == 1 :
			prev_cmd = cmd = X[i]; i += 1
		else :
			cmd = prev_cmd
		if cmd == 'm' :
			a = read_cur()
			print('\draw %s' % p(a))
		elif cmd == 'c' :
			a, b, c = read(), read(), read_cur()
			print('.. controls %s and %s .. %s' % (p(a), p(b), p(c)))
		elif cmd == 'l' :
			a = read_cur()
			print('-- %s' % p(a))
		elif cmd == 'z' :
			print(';')
		else :
			0/0

def g(x) :
	a = BeautifulSoup(x, 'lxml')
	ans = []
	for i in a.select('rect') :
		h = float(i.attrs['height'])
		w = float(i.attrs['width'])
		x = float(i.attrs['x'])
		y = float(i.attrs['y'])
		tr = i.attrs.get('transform')
		co, = re.search('fill:#([0-9a-f]+);', i.attrs['style']).groups()
		color = int(co[0:2], 16), int(co[2:4], 16), int(co[4:6], 16)
		xx = x + w
		yy = y + h
		if tr is not None :
			exp = 'matrix\(([^,]+),([^,]+),([^,]+),([^,]+),([^,]+),([^,]+)\)'
			mat = list(map(float, re.fullmatch(exp, tr).groups()))
			matrix = np.array([[mat[0], mat[2], mat[4]],
								[mat[1], mat[3], mat[5]],
								[0, 0, 1]])
			def transform(x, y) :
				md = matrix.dot(np.array([[x], [y], [1]]))
				return md[0][0], md[1][0]
			xy = transform(x, y)
			xxy = transform(xx, y)
			xyy = transform(x, yy)
			xxyy = transform(xx, yy)
			def mean(a, b) :
				s = (a + b) / 2
				d = abs((a - b) / 2)
				# print(d)	# 0/0
				return s
			if 0 :
				print(x, y)
				print(xx, yy)
				print(xy)
				print(xxy)
				print(xyy)
				print(xxyy)
			xxx = sorted([xy[0], xxy[0], xyy[0], xxyy[0]])
			yyy = sorted([xy[1], xxy[1], xyy[1], xxyy[1]])
			x = mean(xxx[0], xxx[1])
			xx = mean(xxx[2], xxx[3])
			y = mean(yyy[0], yyy[1])
			yy = mean(yyy[2], yyy[3])
		def rnd(x) :
			return round(x / 5, 1) * 5
		ans.append(
			'\\filldraw [color={rgb,255:red,%d; green,%d; blue,%d}]\n' % color+\
			'\t%s rectangle %s;' % (p((rnd(x), rnd(y))), p((rnd(xx), rnd(yy)))))
	for i in reversed(ans) :
		print(i)
#		print(x, y, xx, yy, tr, color)

# f("m 54.697536,20.16266 c -0.479502,-0.0899 -0.882073,-0.173702 -0.894603,-0.186234 -0.01614,-0.01614 1.537356,-8.837361 1.567609,-8.901308 0.004,-0.0085 1.903782,0.330154 1.91559,0.341583 0.0043,0.0041 -0.351867,2.007037 -0.791549,4.450771 -0.439682,2.443735 -0.800263,4.448646 -0.801292,4.45536 -0.0035,0.0227 -0.122021,0.0036 -0.995755,-0.160172 l 0,0 z")

# f("m 54.153405,27.479108 c -0.08056,-0.01027 -0.377179,-0.06158 -0.659148,-0.113991 -3.096766,-0.575626 -5.5713,-3.000914 -6.276113,-6.151211 -0.108815,-0.4864 -0.151208,-2.041513 -0.07301,-2.678493 0.325385,-2.650351 2.107616,-5.022615 4.615841,-6.143985 0.497452,-0.222399 1.441685,-0.519562 1.571482,-0.494565 0.06953,0.01339 0.116632,0.166655 0.245081,0.797359 0.221108,1.085676 0.234817,1.032439 -0.307113,1.192638 -2.095264,0.619375 -3.679526,2.354288 -4.149393,4.543972 -0.121448,0.565954 -0.143853,1.542397 -0.04869,2.121834 0.386054,2.350787 2.139541,4.24826 4.495155,4.864269 0.558722,0.146108 1.782808,0.183379 2.412662,0.07346 0.78671,-0.137295 1.666754,-0.496649 2.364655,-0.965578 0.427868,-0.28749 1.267266,-1.124179 1.548604,-1.543605 0.435943,-0.64992 0.738045,-1.380964 0.922495,-2.232315 0.187192,-0.863992 0.103411,-2.18121 -0.196099,-3.083259 -0.134755,-0.405842 -0.481746,-1.076973 -0.761822,-1.473476 -0.329195,-0.466042 -0.985595,-1.123805 -1.448443,-1.451452 -0.219395,-0.155307 -0.418098,-0.298712 -0.44156,-0.318674 -0.04472,-0.03805 0.419698,-0.959402 0.721046,-1.430485 l 0.175662,-0.274604 0.280952,0.182036 c 0.734761,0.476072 1.570037,1.228968 2.059816,1.85666 0.516011,0.66132 0.995211,1.562257 1.266595,2.381309 0.511224,1.542919 0.51913,3.354742 0.02129,4.882586 -0.520916,1.598699 -1.561923,3.031949 -2.899641,3.992197 -0.822219,0.590209 -2.046963,1.134896 -2.974613,1.322916 -0.780066,0.158108 -1.86219,0.221509 -2.465706,0.144465 z")

g(open('/tmp/logo.svg').read())
\end{comment}

	\end{tikzpicture}
}

\setbeamercolor{background canvas}{bg=}

\newcommand{\PreFirstFrame}{
	\AddToShipoutPictureFG*{
		\AtPageLowerLeft{
			\put(\LenToUnit{0.05\paperwidth},\LenToUnit{0.1\paperheight}){
				\footnotesize
				这个指引文档在
				\href{https://creativecommons.org/licenses/by-sa/3.0/deed.zh}
				{知识共享 署名-相同方式共享 3.0协议} 之条款下提供
			}
			\put(\LenToUnit{0.05\paperwidth},\LenToUnit{0.05\paperheight}){
				\footnotesize
				Available under
				\href{https://creativecommons.org/licenses/by-sa/3.0/}
				{Creative Commons Attribution-ShareAlike License}
			}
			\put(\LenToUnit{0.6\paperwidth},\LenToUnit{0.15\paperheight}){
				\HCCLogoFull
			}
		}
	}
}

\newcommand{\PostFirstFrame}{
	\AddToShipoutPictureBG{
		\AtPageLowerLeft{
			\put(\LenToUnit{0.8\paperwidth},\LenToUnit{0.15\paperheight}){
				\HCCLogoSimp
			}
		}
	}
}

\newcommand{\PreLastFrame}{
	\ClearShipoutPictureBG

	\AddToShipoutPictureFG*{
		\AtPageLowerLeft{
			\put(\LenToUnit{0.6\paperwidth},\LenToUnit{0.15\paperheight}){
				\HCCLogoFull
			}
		}
	}
}

\definecolor{Pink}{HTML}{ff0bff}	% Use: \color{Pink} OR \color[HTML]{ff0bff}
\definecolor{Blue}{HTML}{0000ff}
\definecolor{Green}{HTML}{008a8c}
\definecolor{Brown}{HTML}{a52a2a}

% Note: this C style differs a lot from gedit's
\lstdefinestyle{cstyle}{
	language=c,
	basicstyle=\ttfamily,
	morekeywords={with},
	keywordstyle=\bfseries\color{Brown},
	commentstyle=\color{Blue},
	stringstyle=\color{Pink},
	keywordstyle=[3]\color{Green},
	alsoletter={0123456789.},
	morekeywords=[4]{0,1,2,100,999},
	keywordstyle=[4]\color{Pink},
	upquote=true,
	breaklines=true,
}

\lstdefinestyle{pythonstyle}{
	language=python,
	basicstyle=\ttfamily,
	% frame=single,
	morekeywords={with,yield},
	keywordstyle=\bfseries\color{Brown},
	keywordstyle=[2]\color{Green},
	commentstyle=\color{Blue},
	stringstyle=\color{Pink},
	keywordstyle=[3]\color{Green},
	alsoletter={0123456789.},
	morekeywords=[4]{False,True,
		0,1,2,3,4,5,6,7,8,9,10,11,12,13,15,17,19,16,20,24,27,31,32,33,34,35,38,
		45,56,60,64,81,95,97,99,100,123,243,256,400,512,576,729,999,1000,1024,
		1234,1365,1366,2000,2187,2836,2957,3856,3857,5678,6561,9274,100000,
		1000000,
		0.5,3.14,3.4,
		0x1234,},
	keywordstyle=[4]\color{Pink},
	upquote=true,
	breaklines=true,
	showstringspaces=false,
}

\lstdefinestyle{bashstyle}{
	language=bash,
	basicstyle=\ttfamily,
	% frame=single,
	morekeywords=[2]{
		cp,df,du,find,ll,ls,mkdir,mv,rm,tree,
		cut,diff,grep,head,hexdump,less,man,md5sum,more,nano,sha1sum,tail,vi,wc,
		date,free,ifconfig,lspci,sleep,sync,time,top,uname,
		7za,dnf,history,rpm,su,sudo,wget,yum,
		git,
	},
	keywordstyle=\bfseries\color{Brown},
	keywordstyle=[2]\color{Green},
	commentstyle=\color{Blue},
	stringstyle=\color{Pink},
	keywordstyle=[3]\color{Green},
	alsoletter={-=0123456789.|><&~*?;},
	morekeywords={&,&&,|,||,>,<,>>,<<,2>,&>,;},
	morekeywords=[4]{
		0,1,2,3,4,5,6,7,8,9,10,
		a,b,c,d,e,f,g,h,i,j,k,l,m,n,o,p,q,r,s,t,u,v,w,x,y,z,
		bin,dev,etc,hcc,home,jkl,mnt,opt,root,srv,tmp,usr,
		*,?,.,..,~,-,stdin,stdout,stderr,
		null,zero,random,urandom,sda,sda1,sda2,sda3,sda4,sdb1,sr0,USB,tmpfs,
		passwd,shadow,hosts,fstab,sudoers,locale.conf,resolv.conf,
		.bash_history,mydoc.7z,mydoc,master,
	},
	keywordstyle=[4]\color{Pink},
	upquote=true,
	breaklines=true,
	showstringspaces=false,
}

\lstset{
	tabsize=4,
	columns=fixed,
	extendedchars=false,
}

\newcommand{\inlinePython}{\lstinline[style=pythonstyle]}

\newcommand{\inlineBash}{\lstinline[style=bashstyle]}

\newcommand{\inlineListing}{\lstinline[basicstyle=\ttfamily]}



\begin{document}
\begin{CJK}{UTF8}{gbsn}

\PreFirstFrame
\begin{frame} [fragile]
	\centerline{\fontsize{42}{42}\selectfont Git Talk 1}
\end{frame}
\PostFirstFrame

\begin{frame} [fragile]
	\frametitle{Git是什么}
	\linespread{1.5}
	\begin{itemize}
	\item Git是目前世界上最先进的分布式版本控制系统
		\begin{itemize}
		\item 来自\href{https://www.liaoxuefeng.com/wiki/896043488029600}
						{廖雪峰的官方网站}教程
		\item HCC目前的所有项目，包括幻灯片都在用Git管理
		\end{itemize}
	\end{itemize}
	
	\
	
	\centerline{\includesvg{asserts/git-logo}}
\end{frame}

\begin{frame} [fragile]
	\frametitle{Git设计理念}
	\linespread{1.5}
	\begin{itemize}
	\item Git的项目管理方式
		\begin{itemize}
		\item 工作区
		\item 暂存区
		\item 版本库
		\end{itemize}
	\item 例
	\linespread{0}
	\begin{lstlisting}[style=bashstyle, gobble=4, texcl, escapechar=@]
	shiyiquan			 # 项目目录
	@├─@.git@　@			 # 隐藏目录
	@│　├─@暂存区@　　@		# 临时存储
	@│　└─@版本库@　　@		# 永久存储
	@│　　　├─@98d22f3		  # 一个版本
	@│　　　└─@118532c		  # 另一个版本
	@├─@shiyiquan@　@		 # 工作区文件
	@└─@README.md@　@		 # 可随意更改
	\end{lstlisting}
	\linespread{1.5}
	\end{itemize}
\end{frame}

\begin{frame} [fragile]
	\frametitle{基础操作}
	\linespread{1.25}
	\begin{itemize}
	\item 创建版本库
	\begin{lstlisting}[style=bashstyle, gobble=4, texcl]
	mkdir hcc	# 创建项目目录
	cd hcc		# cd 进去
	git init	# 创建git目录结构
	\end{lstlisting}
	\item 开始工作
	\begin{lstlisting}[style=bashstyle, gobble=4, texcl]
	cat > jkl	# 假设jkl是我们正在编写的文件
	\end{lstlisting}
	\item 提交代码
	\begin{lstlisting}[style=bashstyle, gobble=4, texcl]
	git add jkl		# 工作区$\to$暂存区
	git add -A		# 添加全部文件
	git commit		# 暂存区$\to$版本库
					# 第一次提交需设置用户
					# 和vi操作一样
	\end{lstlisting}
	\end{itemize}
\end{frame}

\begin{frame} [fragile]
	\frametitle{状态查看}
	\linespread{1.25}
	\begin{itemize}
	\item 命令
	\begin{lstlisting}[style=bashstyle, gobble=4, texcl, escapechar=@]
	git status		# 检查状态
	git @diff@		# 对比工作区和暂存区
	git log			# 查看版本库日志
	\end{lstlisting}
	\item 尝试
	\begin{lstlisting}[style=bashstyle, gobble=4, texcl, escapechar=@]
	cat >> tmp		# 输入一些内容
	git add -A
	cat >> tmp		# 再输入一些内容
	git status
	git @diff@
	git log			# 按q退出
	\end{lstlisting}
	\end{itemize}
\end{frame}

\begin{frame} [fragile]
	\frametitle{高级选项}
	\linespread{1.5}
	\begin{lstlisting}[style=bashstyle, gobble=4, texcl, escapechar=@]
	git log
		--oneline	# 摘要模式（左侧是版本号的前几位）
		--graph		# 显示版本关系图（在复杂时后很有用）
	git commit
		-m "@\color{Pink}提交说明@"	# 不进入vi
	\end{lstlisting}
\end{frame}

\begin{frame} [fragile]
	\frametitle{回溯}
	\linespread{1.25}
	\begin{itemize}
	\item 命令
	\begin{lstlisting}[style=bashstyle, gobble=4, texcl, escapechar=@]
	git reset --hard 版本号
	HEAD^		# 当前版本的前一个版本
	HEAD~10		# 前10个版本
	\end{lstlisting}
	\item 回溯后将很难看到更靠后的版本
	\item 查看操作记录
		\begin{itemize}
		\item \inlineBash{git reflog}
		\item 可以看到以后的版本并跳转回去
		\end{itemize}
	\end{itemize}
\end{frame}

\begin{frame} [fragile]
	\frametitle{连接远程版本库}
	\linespread{1.5}
	\begin{itemize}
	\item 通过连接提供git服务的服务器和开发团队同步代码
		\begin{itemize}
		\item \inlineBash{git remote add origin 链接}
		\end{itemize}
	\item 如何获得链接
		\begin{itemize}
		\item 服务商一般会提供链接，复制粘贴即可
		\item 一般链接是有规律的，例如
			\begin{itemize}
			\item \texttt{https://github.com/用户/项目.git}
			\item \texttt{git@github.com:用户/项目.git}
			\end{itemize}
		\item 链接甚至可以是一个本地的git目录
		\item ``\texttt{origin}''是链接的名称，用于快速访问链接
		\end{itemize}
	\end{itemize}
\end{frame}

\begin{frame} [fragile]
	\frametitle{推送和接收}
	\linespread{1.5}
	\begin{itemize}
	\item 推送
		\begin{itemize}
		\item \inlineBash{git push origin master}
		\item 将当前所在分支推送到origin的master分支
		\end{itemize}
	\item 接收
		\begin{itemize}
		\item \inlineBash{git pull origin master}
		\item 将origin的master分支拉到当前分支
		\end{itemize}
	\item ``master''是默认分支名称
	\end{itemize}
\end{frame}

\begin{frame} [fragile]
	\frametitle{分支}
	\linespread{1.25}
	\begin{itemize}
	\item 分支可以代表不同阶段的代码，可以互相合并
	\item 创建
		\begin{itemize}
		\item \inlineBash{git branch dev}
		\item 从当前分支创建dev分支
		\end{itemize}
	\item 跳转
		\begin{itemize}
		\item \inlineBash{git checkout dev}
		\item 切换到dev分支
		\end{itemize}
	\item 合并
		\begin{itemize}
		\item \inlineBash{git merge master}
		\item 将master分支合并到当前分支
		\item 从网络拉取（\inlineBash{git pull}）也可能会合并分支
		\end{itemize}
	\end{itemize}
\end{frame}

\begin{frame} [fragile]
	\frametitle{分支练习}
	\linespread{1.5}
	\begin{itemize}
	\item 从master开始执行
	\begin{lstlisting}[style=bashstyle, gobble=4, texcl, escapechar=@]
	git branch dev
	echo "master ver" >> jkl
	git commit -am "Master"	# -am相当于add后提交
	git checkout dev		# jkl变了
	echo "dev version" >> jkl
	git commit -am "Dev"
	git merge master		# 出现冲突
	\end{lstlisting}
	\end{itemize}
\end{frame}

\begin{frame} [fragile]
	\frametitle{冲突}
	\linespread{1.5}
	\begin{itemize}
	\item 冲突后git会提示冲突的文件，只要按照标记手动合并然后重新提交即可
	\begin{lstlisting}[basicstyle=\ttfamily, gobble=4, texcl, escapechar=@]
	@\color{Pink}<<<<<<< HEAD@			# 当前分支
	dev version
	@\color{Pink}=======@				 # 分割线
	master ver
	@\color{Pink}>>>>>>> master@		  # 合并分支
	\end{lstlisting}
	\end{itemize}
\end{frame}

\begin{frame} [fragile]
	\frametitle{其它功能}
	\linespread{1.5}
	\begin{lstlisting}[style=bashstyle, gobble=4, texcl, escapechar=@]
	git tag					# \href
	{https://www.liaoxuefeng.com/wiki/896043488029600/900788941487552}{标记代码版本}
	git stash				# \href
	{https://www.liaoxuefeng.com/wiki/896043488029600/900785521032192}{隐藏版本}
	git config				# 自定义git
	.gitignore				# \href
	{https://www.liaoxuefeng.com/wiki/896043488029600/900004590234208}{忽略文件}
	git grep				# 全文搜索
	\end{lstlisting}
\end{frame}

\PreLastFrame
\begin{frame}
	\centerline{\fontsize{32}{32}\selectfont 感谢参加此次活动}
\end{frame}

\newpage
\end{CJK}
\end{document}

